\documentclass[10pt]{article}

% Specify the margins and text size.
\setlength{\textwidth}{6.6in}
\setlength{\textheight}{9.5in}
\setlength{\oddsidemargin}{0pt}
\setlength{\evensidemargin}{0pt}
\setlength{\topmargin}{0pt}
\setlength{\hoffset}{-.15in}
\setlength{\voffset}{-1in}

\setlength{\parskip}{5pt}
\setlength{\parindent}{0pt}

% Load some fonts and symbol packages
\usepackage{latexsym}
\usepackage{pifont}       % contains 'star' symbol for counterinsurgency handbook title
\usepackage{yfonts} 
\usepackage{amsmath}
\usepackage{amsfonts}

\usepackage{graphicx}     % actually, this is loaded by pstricks
\usepackage[T1]{fontenc}
\usepackage{ifthen}
\usepackage{pstricks,pst-grad,pst-text,pst-node,multido,pst-plot,calc,pst-3dplot}
%\usepackage[all]{xy}
%\usepackage{animate}

% The hyperref package inserts links into the pdf.
\definecolor{MyLinkColor}{rgb}{.1,.2,1}
\definecolor{MyCiteColor}{rgb}{.1,1,.2}
\definecolor{MyURLColor}{rgb}{.4,.4,.4}
\usepackage[backref=true,pagebackref=false,hyperindex,colorlinks=true,
  linkcolor=MyLinkColor,urlcolor=MyURLColor]{hyperref}


% The tweaklist package is something I found on the web.  It provides a simple interface
% for making changes to spacing used in the itemize and enumerate environments.  Comment
% this out if you don't care to use tweaklist.
\usepackage{tweaklist}
\renewcommand{\itemhook}{\setlength{\parskip}{2pt}\setlength{\parsep}%
{1pt}\setlength{\topsep}{0pt}\setlength{\itemsep}{0pt}}

\newcommand{\U}{\underline{\hspace{5pt}}}

\usepackage{listings}

\begin{document}
\pagestyle{empty}
\lstset{language=R, showspaces=false, showstringspaces=false}

\href{http://www.su.edu}{\includegraphics[height=1.75cm]{sulogo.eps}}
\vspace{-1.69cm}

{\small{\ }\hfill
\begin{tabular}{cl}
& Math 207\\
& Introduction to Statistics\\
\end{tabular}
}
\setlength{\baselineskip}{1.05\baselineskip}

\begin{center}
\textbf{\large  Regression Example}
\end{center}
\newcommand{\Z}{\hphantom{0}}
\newcommand{\VP}{\vphantom{\Large Y}}
\newcommand{\M}{\hphantom{$-$}}

\begin{center}
\begin{tabular}{|rr|ccc||cc|}\hline
    &     &       &       &           & $y$ values    &             \VP\\
    &     &       &       &           & predicted     & regression     \\
$x$ & $y$ & $z_x$ & $z_y$ & $z_x\,z_y$& by regression & error           \\\hline
$1$   & $10$   & $-3/2$\M  & $1/2$  & $-3/4$\M    & $7.85$ & $2.15$ \VP\\
$3$   & $9$   & $-1/2$\M  & $1/4$   & $-1/8$\M    & $7.95$ & $1.05$    \\
$4$   & $8$   & $0$     & $0$       & $0$         & $8.00$ & $0.00$    \\
$4$   & $7$   & $0$     & $-1/4$\M  & $0$         & $8.00$ & $-1.00$\M \\
$5$   & $1$   & $1/2$   & $-7/4$\M  & $-7/8$\M    & $8.05$ & $-7.05$\M \\
$7$   & $13$  & $3/2$   & $5/4$   & $15/8$\Z      & $8.15$ & $4.85$    \\\hline
\end{tabular}
\end{center}

\begin{itemize}
\item  Compute the necessary statistics.
\begin{align*}
\mbox{mean}(x) & = \frac{1}{6}\,(1 + 3 + 4 + 4 + 5 + 7) = 4\\
\mbox{sd}(x)   & = \sqrt{\frac{(1-4)^2 + (3-4)^2 + (4-4)^2 + (4-4)^2 + (5-4)^2 + (7-4)^2}{5}} = 
   \sqrt{\frac{9 + 1 + 0 + 0 + 1 + 9}{5}}\\ & = \sqrt{\frac{20}{5}}=\sqrt{4} 
                =2\\
\mbox{mean}(y) &=\frac{1}{6}\,(10 + 7 + 9 + 8  + 1 + 13) = 8\\
\mbox{sd}(y)   &=\sqrt{\frac{(10-8)^2 + (7-8)^2 + (9-8)^2 + (8-8)^2 + (1-8)^2 + (13-8)^2}{5}}
     = \sqrt{\frac{4 + 1 + 1 + 0 + 49 + 25}{5}} \\ & = \sqrt{\frac{80}{5}}
     = 4
\end{align*}
\item Convert the $x$ values to standard units.  For example,
\[x = 1\mbox{\hspace{10pt}becomes\hspace{10pt}} z_x\frac{1-4}{2} = -3/2
\mbox{\hspace{20pt}and\hspace{20pt}}
x = 3\mbox{\hspace{10pt}becomes\hspace{10pt}} z_x\frac{3-4}{2} = -1/2\]
\item Convert the $y$ values to standard units:
For example,
\[y=10\mbox{\hspace{10pt}becomes\hspace{10pt}} z_y=\frac{10-8}{4} = 1/2
\mbox{\hspace{20pt}and\hspace{20pt}}
y=7\mbox{\hspace{10pt}becomes\hspace{10pt}}  z_y\frac{7-8}{4} = -1/4\]
\item Compute the products $z_x\,z_y$.
\item Compute the correlation coefficient.
\[r = \frac{1}{5}\,\left(-\frac{3}{4} - \frac{1}{8} + 0 + 0 - \frac{7}{8} + \frac{15}{8}\right)
  = \frac{1}{5}\cdot\frac{1}{8} = \frac{1}{40}=0.025\]
%
\item Find the regression line.\vspace{-15pt}
\begin{align*}
\big(y-\mbox{mean}(y)\big) &= r\,\frac{\mbox{sd}_y}{\mbox{sd}_x}\,
   \big(x-\mbox{mean}(x)\big)\\
(y-8) & = \frac{1}{40}\cdot\frac{4}{2}\,(x-4)\\
(y-8) & = \frac{1}{20}\,(x-4)\\
20\,(y-8) &= x-4\\
20\,y-160  & = x - 4\\
20\,y     & = x  + 156\\
y        & = \frac{x + 156}{20}\\
         & = 0.05\,x + 7.8
\end{align*}
%
\item Compute the $y$ values predicted by the regression line.  If $x=1$, then
\[y=(1 + 156)/20= 157/20 = 7.85\]
and if $x=3$, then
 \[y=(3 + 156)/20 = 159/20 = 7.95.\]
%
\item Compute the regression errors $(\mbox{error}=y - \mbox{predicted $y$})$.  
For example,
\[\mbox{if $x=1$, then $\mbox{error}=10-7.85=2.15$\hspace{20pt} and\hspace{20pt} if $x=3$, then
  $\mbox{error}=9-7.95=1.05$.}\]
\item Compute the RMS size of the errors.
\[\mbox{RMS}_{\mbox{\scriptsize reg}} = \sqrt{(2.15^2 + 1.05^2 + 0^2 + (-1)^2 + (-7.05)^2 + 4.85^2)/6} 
  = \sqrt{(79.95)/6}
  = \sqrt{13.325} \approx 3.65.\]
\end{itemize}
\bigskip

Here is how to get \texttt{R} to do all this:

\begin{tabular}{ll}
\texttt{x <- c(1,  3, 4, 4, 5, 7)} & define the $x$ values\\
\texttt{y <- c(10, 9, 8, 7, 1, 13)} & define the $y$ values\\
\texttt{mean(x)}                    & We can actually skip from here $\dots$\\
\texttt{sd(x)}   &\\
\texttt{mean(y)} & \\
\texttt{sd(y)}\\
\texttt{zx <- (x - mean(x)) / sd(x)} & \\
\texttt{zy <- (y - mean(y)) / sd(y)} & \\
\texttt{r <-  mean(zx * zy)} & \\
\\
\texttt{slope <- r * sd(y) / sd(x)} & \\
\texttt{yIntercept <- mean(y) - slope * mean(x)}  &  
         \hspace{30pt} $\dots$ to the next line\\[5pt]
\texttt{linearModel <- lm(y $\sim$ x)} \\
\texttt{summary(linearModel)} & read the slope and $y$ intersecpt from this\\
\texttt{fitted.values(linearModel)} & this gives the $y$ values predicted by regression\\
\texttt{RegressionErrors <- residuals(linearModel)} & \\
\texttt{SDrms <- sqrt(mean(RegessionErrors**2))}
\end{tabular}

\begin{center}
\begin{pspicture}(0,0)(7,8)
\psset{yunit=0.6,dotsize=0.15}
\psdot(1,10)
\psdot(3,9)
\psdot(4,8)
\psdot(4,7)
\psdot(5,1)
\psdot(7,13)
\psline[linewidth=0.02](0,7.8)(7,8.15)
\rput(3.5,-1.7){$x$}\rput(-1,6.5){$y$}
\psaxes(0,0)(7,13)
\end{pspicture}
\end{center}
\vfill
\eject

\end{document}

